\documentclass{article}
\usepackage{geometry}  
\usepackage{amsmath}  
\usepackage{graphicx} 
\usepackage{svg}
\usepackage[T1]{fontenc}
\usepackage[backend=biber,style=numeric]{biblatex}
\usepackage{url} % Obsługa linków
\usepackage{hyperref}


%----------------------------------------------------------------------------------------
%	REPORT INFORMATION
%----------------------------------------------------------------------------------------

\title{TRA - Techniki Realizacji Algorytmów Cyfrowego Przetwarzania Sygnałów  \\ Projekt - specyfikacja wstępna \\} % Report title

\author{Anna \textsc{Dzieżyk}, indeks: 318506 \\ }% Author name(s), add additional authors like: '\& James \textsc{Smith}'

\date{\today} % Date of the report

%----------------------------------------------------------------------------------------

\begin{document}

\maketitle % Insert the title, author and date using the information specified above
\section{Cel projektu}
Celem projektu będzie stworzenie stopnia wejściowego dla przykładowo wzmacniacza klasy D użytego w kolumnie głośnikowej -  
ten projektowany stopień wejściowy nie będzie komparatorem tylko modulatorem sigma delta. Następnie zostanie porównany typowy model 
wzmacniacza klasy D z komparatorem na wejśćciu, do zaprojektowanego wzmacniacza z modulatorem sigma-delta na wejściu.
\section{Wstępny sposób realizacji}

Stworzenie modelu wzmacnacza klasy D, z komparatorem jako stopniem wejściowym, podanie spróbkowanego sygnału audio na wejście wzmacniacza, następnie stworzenie modelu 
wzmacniacza klasy D z modulatorem Sigma-Delta (od pierwszego do trzeciego rzędu) jako stopniem wejściowym i finalnie porównanie parametrów wyjściowych sygnałów obserwowanych na wyściu wszystkich zaprojektowanych wzmacniaczy.
\section{Potrzebne środowiska}
Matlab, Simulink;
\section{Bibliografia}
1. \url{https://pl.wikipedia.org/wiki/Wzmacniacz_impulsowy} \\
2. \url{https://www.digikey.pl/pl/articles/the-basics-of-anti-aliasing-low-pass-filters}
3. \url{https://esezam.okno.pw.edu.pl/mod/book/view.php?id=33&chapterid=648}
\end{document}